\chapter{44}



Mehrere Polizeisirenen zerrissen kurz vor halb fünf die Nacht. Die
Ambulanz war zuerst da. Kurze Zeit später kamen ein Polizeiwagen und die
Feuerwehr angerast.

Ein Mann war gehetzt auf der Polizeistation aufgetaucht. Rund zwanzig
Minuten zuvor.

Mit einem kraftvollen Ruck hatte er die Glastür des Grossraumbüros an
die Innenwand geschlagen. Hatte sein T-Shirt am Halsausschnitt bis über
seine Nase hochgezogen und ungebremst einen Schwall Worte hineingeworfen
und der Polizistin mitten ins Gesicht. Die Frau sass alleine im Raum.
Hielt sich duckend am Rand des Schreibtisches fest und starrte ihn
verdattert an.

«Beim \emph{«Roccia del Culto»}, südlich der Felswand, fast unten, liegt
ein Mensch auf einem Baum, hab’ nen Arm gesehen, weiss nicht, auf der
Krone oben, bewegungslos, keine Ahnung, wie und warum.»

Noch bevor sie nach Details fragen konnte, war er verschwunden. Ihre
Kollegen kamen herbeigeschossen. Sie hatten einen Knall und eine fremde,
aufgeregte Stimme gehört. Sie befürchteten einen Überfall.

Ausser der Polizistin hatte den Mann niemand gesehen.

Sie erinnerte sich anschliessend vorwiegend an die auffallend hohe
Stimme und an seinen Blick. Mittleren Alters, mittlerer Grösse,
vielleicht um die 40 war er, schätzte sie, mit ungewöhnlichen Augen.
Dass ihr war, als sehe er durch sie hindurch.

Niemand konnte sich erst erklären, wie er es ungehindert in die
Büroräume geschafft hatte. Am Eingang vorbei, wo der diensthabende
Kollege sass, oder, hätte sitzen sollen.

\sectionbreak

Commissario Alessandro Mantovani und ein weiterer Polizist blickten
gebannt dem Baumstamm entlang nach oben. Sie ahnten Schlimmes. Im Korb
der Feuerwehrleiter sollten zwei Sanitäter hochgefahren werden.

Die Bergung war kompliziert. Der Feuerwehrwagen konnte nicht direkt
unter die Akazie gefahren werden. Der Baum lag etwas erhöht oberhalb
eines schmalen Pfades. Inmitten einer ganz kleinen Wiese. Sie befand
sich als grössere Einbuchtung im untersten Teil der Felswand. Der Fahrer
liess die Feuerwehrleiter senkrecht, anschliessend rechtwinklig und bis
zum Anschlag langsam hochfahren. Schliesslich konnte er die seitlichen
Abstützungen ausfahren und das schwere Gefährt sichern. Dann übernahm
der Maschinist die Steuerung am Drehkreuz. Sie hatten Glück. Es gelang,
den Leiterkorb so nahe an den auf dem Baum liegenden Menschen zu
steuern, dass sich einer der Sanitäter direkt neben ihm befand.

«Eine junge Frau,» rief dieser aufgeregt in sein Funkgerät nach unten,
während er ihren Puls fühlte, «sie lebt, bewusstlos, aber sie lebt.»

Mantovani und die andern Anwesenden atmeten auf.
